\documentclass[10pt]{article}
\usepackage[T1]{fontenc}
\usepackage[utf8]{inputenc}
\usepackage[french]{babel}
\usepackage{amsmath,amssymb,mathrsfs,mathpazo}
\usepackage{geometry}
\geometry{a4paper,margin=2cm}
\usepackage{graphicx,array,tabularx,tabularray,colortbl,multirow,multicol,booktabs,enumitem,pifont,ulem,fancybox,diagbox,cancel}
\usepackage{pstricks,pstricks-add,pst-plot,pst-tree,pst-node,pst-text,pst-func,pst-math,pst-bspline,pst-3dplot,pst-eucl,pst-solides3d}
\let\pscalculate\undefined
\usepackage{pst-bezier}
\usepackage{tikz,pgfplots,tkz-tab}
\usetikzlibrary{arrows,arrows.meta,patterns,calc,positioning,trees,decorations.markings,shapes.geometric,shapes.misc,shapes.arrows,intersections,decorations.pathmorphing,decorations.pathreplacing,backgrounds,shadows,babel}
\pgfplotsset{compat=1.18}
\usepgfplotslibrary{fillbetween}
\usepackage[np]{numprint}
\usepackage{listings,makecell,pas-tableur,twoopt}
\let\barre\undefined
\usepackage[pstricks]{tablvar}
\definecolor{darkgreen}{RGB}{0,100,0}
\providecommand{\red}{\color{red}}
\providecommand{\blue}{\color{blue}}
\providecommand{\green}{\color{green}}
\providecommand{\black}{\color{black}}
\providecommand{\white}{\color{white}}
\providecommand{\gray}{\color{gray}}
\providecommand{\cyan}{\color{cyan}}
\providecommand{\magenta}{\color{magenta}}
\providecommand{\yellow}{\color{yellow}}
\providecommand{\violet}{\color{violet}}
\providecommand{\orange}{\color{orange}}
\providecommand{\euro}{\texteuro}
\providecommand{\EUR}{\texteuro}
\providecommand{\ieme}{\textsuperscript{e}}
\providecommand{\degres}{\ensuremath{^\circ}}
\providecommand{\degree}{\ensuremath{^\circ}}
\providecommand{\vv}[1]{\overrightarrow{#1}}
\providecommand{\Coord}[1]{#1}
\providecommand{\decosix}{}
\providecommand{\starredbullet}{\ensuremath{\bullet}}
\pagestyle{empty}
\setlength{\parindent}{0pt}
\usepackage{adjustbox}\DeclareUnicodeCharacter{2236}{:}\DeclareUnicodeCharacter{2190}{\ensuremath{\leftarrow}}
\providecommand{\R}{\mathbb{R}}
\providecommand{\N}{\mathbb{N}}
\providecommand{\Z}{\mathbb{Z}}
\providecommand{\Q}{\mathbb{Q}}
\providecommand{\D}{\mathbb{D}}
\providecommand{\C}{\mathbb{C}}
\providecommand{\e}{{\mathrm{e}}}
\providecommand{\vect}[1]{\overrightarrow{#1}}
\providecommand{\vectt}[1]{\overrightarrow{\mathrm{#1}}}
\providecommand{\vvt}[1]{\overrightarrow{\mathrm{#1}}}
\providecommand{\np}[2][]{{#2~#1}}
\providecommand{\vv}[1]{{\overrightarrow{#1}}}
\renewcommand{\d}{{\mathrm{d}}}
\providecommand{\degres}{{^\circ}}
\providecommand{\degre}{{^\circ}}
\providecommand{\blue}{\color{blue}}
\providecommand{\red}{\color{red}}
\providecommand{\enskip}{\quad}
\providecommand{\llbracket}{\lbrack\!\lbrack}
\providecommand{\rrbracket}{\rbrack\!\rbrack}
\providecommand{\euro}{{\text{€}}}
\providecommand{\no}{{\text{n°}}}
\providecommand{\interval}[2]{{[#1;#2]}}
\providecommand{\intervFO}[2]{{[#1;#2[}}
\providecommand{\dotfill}{\ldots}

\providecommand{\ds}{\displaystyle}
\providecommand{\textding}[1]{#1}
\providecommand{\Oij}{\(\mathrm{O};\vec{\imath},\vec{\jmath}\)}
\providecommand{\Oijk}{\(\mathrm{O};\vec{\imath},\vec{\jmath},\vec{k}\)}
\providecommand{\Ouv}{\(\mathrm{O};\vec{u},\vec{v}\)}
\providecommand{\intervFF}[2]{\ensuremath{[#1;#2]}}
\renewcommand{\euro}{\ifmmode\text{\texteuro}\else\texteuro\fi}
\DeclareUnicodeCharacter{00B0}{\ifmmode^{\circ}\else\textdegree\fi}
\DeclareUnicodeCharacter{20AC}{\ifmmode\text{\texteuro}\else\texteuro\fi}
% Énoncés : Éducation nationale. Transcription et corrigés : APMEP, Denis Vergès et contributeurs. Découpage et classement : CoopMaths, Sébastien Lozano, Éric Elter et contributeurs. Adaptation de la présentation : Maths Explorer.
% CC BY-SA : https://coopmaths.fr/www/about/
% https://coopmaths.fr/alea/static/bac/2024/tex/bac_2024_09_sujet1_metropole_3_cor.tex
\begin{document}
\noindent\textit{APMEP / CoopMaths. CC BY-SA. Présentation adaptée par Maths Explorer.}\par\medskip
\smallskip



\begin{enumerate}
\item On considère une suite $\left(t_n\right)$ vérifiant la relation de récurrence : \\
pour tout entier naturel $n$, $t_{n+1}= - 0,8t_n + 18.$

	\textbf{Affirmation 1} : La suite $\left(w_n\right)$ définie pour tout entier naturel $n$ par $w_n = t_n - 10$ est géométrique.

\smallskip
	
Pour tout $n$: 
$\begin{aligned}[t]
w_{n  +1} & =t_{n+1}-10 = -0,8 t_n +18 - 10 = -0,8t_n +8 = -0,8\left ( t_n-10\right)\\
&  = -0,8 w_n
\end{aligned}$

Donc la suite $(w_n)$ est géométrique de raison $-0,8$.
	
\hfill{\blue\bf Affirmation 1 vraie}	
	
\item Soit la suite $\left(S_n\right)$ qui vérifie pour tout $n\in \N^*$: $3n - 4 \leqslant S_n \leqslant 3n + 4.$

	La suite $\left(u_n\right)$ est définie, pour tout entier naturel $n$ non nul, par : $u_n=\dfrac{S_n}{n}$.
	
	\textbf{Affirmation 2} : La suite $\left(u_n\right)$ converge.
	
	\smallskip

$n>0$ donc \\
$3n - 4 \leqslant S_n \leqslant 3n + 4.$ entraîne $\dfrac{3n - 4}{n} \leqslant \dfrac{S_n}{n} \leqslant \dfrac{3n + 4}{n}$ c'est-à-dire
$3-\dfrac{4}{n} \leqslant u_n \leqslant 3+\dfrac{4}{n}$.

\begin{list}{\textbullet}{}
\item $\ds\lim_{n\to +\infty} \dfrac{4}{n}=0$ donc $\ds\lim_{n\to +\infty} 3-\dfrac{4}{n}=3$
\item $\ds\lim_{n\to +\infty} \dfrac{4}{n}=0$ donc $\ds\lim_{n\to +\infty} 3+\dfrac{4}{n}=3$
\item Pour tout $n>0$: $3-\dfrac{4}{n} \leqslant u_n \leqslant 3+\dfrac{4}{n}$
\end{list}

Donc, d'après le théorème des gendarmes, on peut déduire que la suite $(u_n)$ est convergente et que $\ds\lim_{n\to +\infty} u_n=3$.

\hfill{\blue\bf Affirmation 2 vraie}	

\item Soit la suite $\left(v_n\right)$ définie par :
$ v_1=2 \text{ et pour tout entier naturel } n\geqslant1,~v_{n+1}=2-\dfrac{1}{v_n}. $
	
\textbf{Affirmation 3} : Pour tout entier naturel $n \geqslant 1$, $v_{n}=\dfrac{n+1}{n}$.
	
\smallskip
	
On va démontrer cette propriété par récurrence.

\begin{list}{\textbullet}{}
\item  \textbf{Initialisation}

Pour $n=1$: $v_1=2$ et $\dfrac{n+1}{n}=\dfrac{1+1}{1}=2$. Donc la propriété est vraie au rang 1.

\item \textbf{Hérédité}

On suppose la propriété vraie au rang $n$, c'est-à-dire $v_n=\dfrac{n+1}{n}$ (hypothèse de récurrence).

$\begin{aligned}
v_{n+1} & = 2-\dfrac{1}{v_n} = 2 - \dfrac{1}{\frac{n+1}{n}} = 2- \dfrac{n}{n+1} = \dfrac{2(n+1)-n}{n+1} = \dfrac{2n+2-n}{n+1} = \dfrac{n+2}{n+1}\\
&  = \dfrac{(n+1)+1}{n+1}
\end{aligned}$

Donc la propriété est vraie au rang $n+1$.

\item \textbf{Conclusion}

La propriété est vraie au rang 1, et elle est héréditaire pour tout $n\geqslant 1$. Donc, d'après le principe de récurrence, la propriété est vraie pour tout $n\geqslant 1$.
\end{list}	
	
\hfill{\blue\bf Affirmation 3 vraie}		
	
\item On considère la suite $\left(u_n\right)$ définie pour tout entier naturel $n$ par $u_n = \e^n-n$.
	
	\textbf{Affirmation 4} : La suite $\left(u_n\right)$ converge.
	
	\smallskip
	
Soit $f$ la fonction définie sur $\R$ par: $f(x)= \e^{x}-x$. Donc $u_n=f(n)$ pour tout $n$.

Pour $x>0$, $f(x)=x \left ( \dfrac{\e^{x}}{x} - 1 \right )$.

On sait que $\ds\lim_{x\to +\infty} \dfrac{\e^{x}}{x} = +\infty$ donc $\ds\lim_{x\to +\infty} \dfrac{\e^{x}}{x} -1 = +\infty$ et donc $\ds\lim_{x\to +\infty} x\left ( \dfrac{\e^{x}}{x}-1 \right ) = +\infty$.

On a alors $\ds\lim_{x\to +\infty} f(x) = +\infty$n ce qui entraîne que $\ds\lim_{n\to +\infty} u_n = +\infty$.

\hfill{\blue\bf Affirmation 4 fausse}		
	
\item On considère la suite $\left(u_n\right)$ définie à l'aide du script écrit ci-dessous en langage \textsf{Python}, qui renvoie la valeur de $u_n$.
	
\texttt{\begin{center}
\begin{tabular}{|l|}\hline
\textbf{def} u(n) :\\
\qquad valeur = 2\\
\qquad \textbf{for} k \textbf{in range}(n) :\\
\qquad \qquad valeur = 0.5 * (valeur + 2/valeur)\\
\qquad \textbf{return} valeur\\ \hline
\end{tabular}
\end{center}
}

On admet que $\left(u_n\right)$ est décroissante et vérifie pour tout $n$ :
$\sqrt{2} \leqslant u_n \leqslant 2.$

\textbf{Affirmation 5} : La suite $\left(u_n\right)$ converge vers $\sqrt{2}$.

\smallskip

D'après le script Python, $(u_n)$ est définie par:
$\left \{
\begin{array}{l !{=} l}
u_0 & 2\\
u_{n+1} & 0,5 \left (u_n + \frac{2}{u_n} \right) \text{ pour tout }n
\end{array}
\right .$

\begin{list}{\textbullet}{}
\item La suite $(u_n)$ est décroissante.
\item Pour tout $n$, on a: $\sqrt{2} \leqslant u_n \leqslant 2$ donc la suite est minorée par $\sqrt{2}$.
\item D'après le théorème de la convergence monotone, on peut dire que la suite $(u_n)$ converge vers un réel $\ell$ tel que $\ell \geqslant \sqrt{2}$.
\end{list}

$\ds\lim_{n\to + \infty} u_n = \ell$ et $\ds\lim_{n\to + \infty} u_{n+1} = \ell$

De l'égalité $u_{n+1} = 0,5 \left ( u_n + \dfrac{2}{u_n} \right )$, on déduit:
$\ell = 0,5 \left ( \ell + \dfrac{2}{\ell} \right )$. 

On résout cette équation.

$\ell = 0,5 \left ( \ell + \dfrac{2}{\ell}\right )
\iff
2\ell = \ell + \dfrac{2}{\ell}
\iff
\ell=\dfrac{2}{\ell}
\iff
\ell^2=2
\iff
\ell=\sqrt{2} \text{ ou } \ell=-\sqrt{2}$

On a vu que $\ell\geqslant \sqrt{2}$ donc $\ell=\sqrt{2}$.

\hfill{\blue\bf Affirmation 5 vraie}		

\end{enumerate}

\end{document}
